\documentclass{article}
\usepackage{amsmath}
\usepackage{amssymb}
\usepackage{hyperref}
\usepackage{braket}

\title{Lecture Notes on Quantum Support Vector Machines}
\author{MHJ}
\date{\today}

\begin{document}

\maketitle

\tableofcontents
\newpage

\section{Introduction to Support Vector Machines}

\subsection{Basic Concepts}
Support Vector Machines (SVM) are supervised learning algorithms used for classification tasks. The main goal of SVM is to find the optimal separating hyperplane (in high-dimensional space) that provides a maximum margin between classes.

\subsection{Mathematical Formulation}
For a dataset \((\mathbf{x}_i, y_i)\) where \(\mathbf{x}_i \in \mathbb{R}^n\) and \(y_i \in \{-1, 1\}\), the decision boundary is defined as:

\[ f(\mathbf{x}) = \mathbf{w} \cdot \mathbf{x} + b = 0 \]

The goal is optimizing:

\[
\min_{\mathbf{w}, b} \frac{1}{2} ||\mathbf{w}||^2
\]

Subject to:

\[
y_i (\mathbf{w} \cdot \mathbf{x}_i + b) \geq 1, \quad \forall i
\]

\section{Quantum Computing Basics}

\subsection{Quantum Bits and Gates}
A qubit can encode a 0 and 1 simultaneously, represented as \(\ket{\psi} = \alpha\ket{0} + \beta\ket{1}\), where \(|\alpha|^2 + |\beta|^2 = 1\).

\subsection{Entanglement and Superposition}
Entanglement is a key resource for quantum computing, enabling qubits to exhibit correlations that are not possible classically. Superposition allows a quantum system to be in multiple states simultaneously, leading to exponential parallelism.

\section{Quantum Support Vector Machines}

\subsection{Motivation}
QSVM leverages quantum computations such as quantum phase estimation and quantum matrix inversion to enhance SVM algorithms, particularly in efficiently handling large datasets and complex kernels.

\subsection{Quantum Kernel Estimation}
In classical SVM, kernels help with non-linear data separations. Quantum computers can speed up the computation of complex kernel evaluations by efficiently simulating an inner product in an exponentially large Hilbert space.

\subsubsection{Quantum Kernel Trick}
Similar to classical SVM, QSVM utilizes a \textit{quantum-enhanced kernel}:

\[
K(\mathbf{x}, \mathbf{y}) = |\braket{\phi(\mathbf{x}) | \phi(\mathbf{y})}|^2
\]

Here, \(|\phi(\mathbf{x})\rangle\) is the quantum state encoding of the classical data \(\mathbf{x}\).

\section{Quantum Advantage in SVM}

\subsection{Quantum Speedup}
Quantum algorithms like HHL (Harrow, Hassidim, and Lloyd) algorithm for solving linear equations provides polynomial speedup, exploiting quantum parallelism and entanglement.

\subsection{Practical Considerations}
Actual implementation challenges include qubit coherence times, error rates, and noise management, alongside classical preprocessing strategies to leverage quantum-enhanced procedures efficiently.

\section{Applications}
Quantum Support Vector Machines present vast potential in finance for fraud detection, in healthcare for diagnosing conditions from large biological datasets, and broadly in any area requiring rapid classification and pattern recognition above classical limits.

\section{Conclusion and Future Work}
QSVM embodies promising potential advancements in quantum machine learning. The road forward involves demonstrating tangible quantum advantages on existing quantum hardware, advancing error correction techniques, and developing larger qubit systems.

\end{document}
